\medskip
\noindent
13.4. \underline{Interpretation in terms of Deligne's ``mixed Hodge
structures''.}

Let us first return to the dictionary of 13.1.4, which we shall
reformulate in a slightly different way, better suited to introducing
the ``Hodge structures'' relevant to the present situation.  For this
purpose, let $\underline R$ be the sheaf of sections of $R$ over $S$, and
introduce the $\underline{\mathcal O}_S$-module

\begin{equation}\tag{13.4.1}
 \mathcal E(A)=\mathcal E=
 \underline R\otimes_{\mathbb Z}\underline{\mathcal O}_S\qquad .
\end{equation}

\noindent
It is locally free of finite type when $R$ is locally constant.  (Note
that $\mathcal E$ need not be coherent when $R$ is not locally
constant.)  The homomorphism (13.1.3)
$\underline R\longrightarrow\underline E$ then defines a canonical
homomorphism

\begin{equation}\tag{13.4.2}
 \mathcal E=\underline R\otimes_{\mathbb Z}
 \underline{\mathcal O}_S\longrightarrow\underline E\qquad .
\end{equation}

\noindent
We denote its kernel by $\underline F(A)$, or simply by $\underline F$;
thus there is an exact sequence

\begin{equation}\tag{13.4.3}
 0\longrightarrow\underline F\longrightarrow\mathcal E
 \longrightarrow\underline E\qquad .
\end{equation}

\noindent
The most interesting case is that in which (13.4.2) is an epimorphism,
that is, in which $R_s$ spans the complex vector space $E_s$ for every
$s\in S$.  We then have an exact sequence

\begin{equation}\tag{13.4.4}
 0\longrightarrow\underline F\longrightarrow\mathcal E
 \longrightarrow\underline E\longrightarrow0\qquad .
\end{equation}

\noindent
If $R$ is locally constant, this is a short exact sequence of locally
free modules over $S$.  Knowledge of the pair $(E,R)$ is then equivalent
to knowledge of the pair $(R,\underline F)$, where $R$ is a locally
constant analytic group over $S$ whose fibres are free $\mathbb Z$-modules
of finite type, and $\underline F$ is a locally direct-summand submodule
of $\mathcal E=\underline R\otimes_{\mathbb Z}
\underline{\mathcal O}_S$, subject in addition to

\begin{equation}\tag{13.4.5}
 \underline F\cap\underline R=0\qquad .
\end{equation}

\noindent
From the standpoint of Hodge structures, $\mathcal E$ should be regarded
as filtered by the ``\underline{Hodge filtration}''

\begin{equation}\tag{13.4.6}
 \operatorname{Fil}^{1}(\mathcal E)=\mathcal E,
 \qquad \operatorname{Fil}^{0}(\mathcal E)=\underline F,
 \qquad \operatorname{Fil}^{-1}(\mathcal E)=0\qquad .
\end{equation}

\noindent
Notice that the exact sequence (13.4.3) still depends functorially on the
analytic group $A$ over $S$, and commutes with arbitrary base change.

Returning now to the conditions of 13.3, put

\begin{equation}\tag{13.4.7}
 \mathcal E^f=\underline R^f\otimes_{\mathbb Z}
 \underline{\mathcal O}_S,
 \qquad
 \mathcal E^t=\underline R^t\otimes_{\mathbb Z}
 \underline{\mathcal O}_S\qquad .
\end{equation}

\noindent
These modules are defined on a suitable neighbourhood of $s$, which we
may suppose to be all of $S$, and define a filtration of $\mathcal E$,
called the \underline{weight filtration},

\begin{equation}\tag{13.4.8}
 \mathcal E\supset\mathcal E^f\supset\mathcal E^t\supset0\qquad .
\end{equation}

\noindent
Here $\mathcal E^t$, $\mathcal E^f$, and $\mathcal E$ are to be regarded
as having weights $-2$, $-1$, and $0$, respectively; thus this is an
increasing filtration.  Like $R^t$, the module $\mathcal E^t$ depends on
the choice of a subtorus $T_o$ of $A_o=A_s$.  Since $R^f$ and $R^t$ are
locally constant, \underline{$\mathcal E^f$ and $\mathcal E^t$ are
locally free}; they may also be interpreted as

\begin{equation}\tag{13.4.9}
 \mathcal E^f=\mathcal E(A^\natural),
 \qquad \mathcal E^t=\mathcal E(T)\qquad .
\end{equation}

\noindent
Applying the functoriality of (13.4.3), as $A$ varies, to the
homomorphisms

\begin{equation*}
 T\longrightarrow A^\natural\longrightarrow A
\end{equation*}

\noindent
shows at once that $\underline F(A^\natural)$ and $\underline F(T)$
identify respectively with

\begin{equation}\tag{13.4.10}
 \underline F(A^\natural)\simeq
 \underline F^f\stackrel{\mathrm{def}}{=}
 \underline F\cap\mathcal E^f,
 \qquad
 \underline F(T)=\underline F^t=
 \underline F\cap\mathcal E^t=0\qquad .
\end{equation}

\noindent
The last equality follows immediately from the fact that
$\underline F(T)=0$ for a torus $T$: in this case the exact sequence of
type (13.4.3) is
$0\longrightarrow0\longrightarrow\mathcal E(T)
\longrightarrow\underline{\operatorname{Lie}}(T)\longrightarrow0$.
At the same time we see that the composite homomorphism

\begin{equation*}
 \mathcal E^t\longrightarrow\mathcal E\longrightarrow\underline E
\end{equation*}

\noindent
induces an isomorphism

\begin{equation}\tag{13.4.11}
 \mathcal E^t\stackrel{\sim}{\longrightarrow}\underline E^t
 \simeq\underline{\operatorname{Lie}}(T)\qquad .
\end{equation}

\noindent
\underline{Suppose now that $R_s$ spans the complex vector space $E_s$},
as happens in particular in the favourable case 13.3.16, where $A_s$ is
an extension of an abeloid group by a complex torus $T_o$.  After
restricting $S$ to a suitable neighbourhood of $s$, the canonical
homomorphism of type (13.4.2), with $A$ replaced by $A^\natural$,

\begin{equation*}
 \mathcal E^f=\mathcal E(A^\natural)
 \longrightarrow\underline{\operatorname{Lie}}(A^\natural)
 \simeq\underline E
\end{equation*}

\noindent
is an epimorphism; a fortiori, so is (13.4.2).  We therefore have the
exact sequences (13.4.4) and

\begin{equation}\tag{13.4.12}
 0\longrightarrow\underline F^f\longrightarrow\mathcal E^f
 \longrightarrow\underline E\longrightarrow0\qquad .
\end{equation}

\noindent
In particular,

\begin{equation}\tag{13.4.13}
 \underline F+\mathcal E^f=\mathcal E\qquad ,
\end{equation}

\noindent
since $\mathcal E/\underline F\simeq\underline E$.  Moreover, since
$B=A^\natural/T$ is a quotient of $A^\natural$,
\footnote{The printed source says here that $B$ is a quotient of $A$;
the displayed definition $B=A^\natural/T$ gives the intended group.}
it satisfies the analogous surjectivity condition, and hence gives an
exact sequence of type (13.4.4):

\begin{equation}\tag{13.4.14}
 0\longrightarrow\underline F(B)\longrightarrow\mathcal E(B)
 \longrightarrow\underline{\operatorname{Lie}}(B)
 \longrightarrow0\qquad .
\end{equation}

\noindent
Its interpretation in terms of the modules already introduced is
immediate.  Indeed, by (13.3.12),

\begin{equation}\tag{13.4.15}
 \left\{
 \begin{array}{l}
  \mathcal E(B)\simeq
  (\underline R^f/\underline R^t)\otimes_{\mathbb Z}
  \underline{\mathcal O}_S
  \simeq\mathcal E^f/\mathcal E^t,\\[2ex]
  \underline{\operatorname{Lie}}(B)
  \simeq\underline E/\underline E^t.
 \end{array}
 \right.
\end{equation}

\noindent
It follows from (13.4.11) and the relation
$\underline F\cap\mathcal E^t=0$ in (13.4.10) that
$\underline F(B)=\ker(\mathcal E^f/\mathcal E^t
\longrightarrow\underline E/\underline E^t)$ is isomorphic to
$\underline F^f$ through the homomorphism

\begin{equation}\tag{13.4.16}
 \underline F^f\simeq\underline F(A^\natural)
 \stackrel{\sim}{\longrightarrow}\underline F(B)\qquad .
\end{equation}

\noindent
Let us summarize the properties just obtained for the weight filtration
(13.4.8), and its relation with the Hodge filtration (13.4.6) of
$\mathcal E$ defined by the submodule $\underline F$, in the favourable
case 13.3.16 where $A_s$ is an extension of an abeloid analytic group
$B_o$ by a torus $T_o$.  We use this torus to define $R^t$, so that $B$
is an abeloid analytic group over $S$.

\medskip
\noindent
\underline{Proposition 13.4.17.} \underline{Under the preceding
conditions, the Hodge and weight filtrations on}
$\mathcal E=\underline R\otimes_{\mathbb Z}\underline{\mathcal O}_S$,

\begin{equation}\tag{13.4.17.1}
 \mathcal E\supset\underline F,
 \qquad
 \mathcal E\supset\mathcal E^f\supset\mathcal E^t\supset0,
\end{equation}

\noindent
\underline{satisfy the following conditions:}

\smallskip
\noindent
a) \underline{The weight filtration is obtained, by tensoring with
$\underline{\mathcal O}_S$ over $\mathbb Z$, from the analogous
filtration on $\underline R$}

\begin{equation*}
 \underline R\supset\underline R^f\supset\underline R^t\supset0,
\end{equation*}

\noindent
\underline{where $\underline R^f$ and $\underline R^t$ are locally
constant analytic groups over $S$, with values free $\mathbb Z$-modules
of finite type}, namely
$R_o^f=R_s=\pi_1(A_s)$ and $R_o^t=\pi_1(T_o)$.

\smallskip
\noindent
b) \underline{One has the relations}

\begin{equation}\tag{*}
 \underline F\cap\mathcal E^t=0,
 \qquad \underline F+\mathcal E^f=\mathcal E\qquad .
\end{equation}

\noindent
\underline{In other words, the filtration on $\mathcal E^t$ induced by
the Hodge filtration}, with the degree conventions of (13.4.6),
\underline{is concentrated in degree $-1$}, whereas the filtration on
$\mathcal E/\mathcal E^f$ induced by the Hodge filtration is
\underline{concentrated in degree $0$}.

\smallskip
\noindent
c) \underline{Consider on}

\begin{equation*}
 \mathcal E^{\mathrm{ab}}
 \stackrel{\mathrm{def}}{=}\mathcal E^f/\mathcal E^t
 \simeq\underline R^{\mathrm{ab}}\otimes_{\mathbb Z}
 \underline{\mathcal O}_S,
 \qquad
 \underline R^{\mathrm{ab}}=\underline R^f/\underline R^t
 \quad\text{(cf. a))},
\end{equation*}

\noindent
\underline{the filtration induced by the Hodge filtration on
$\mathcal E$}.  It is therefore defined by the submodule
$\underline F^{\mathrm{ab}}$ of $\mathcal E^{\mathrm{ab}}$ which is the
image of

\begin{equation*}
 \underline F^f\stackrel{\mathrm{def}}{=}
 \underline F\cap\mathcal E^f\longrightarrow\mathcal E^{\mathrm{ab}}.
\end{equation*}

\noindent
The arrow is a monomorphism by the first relation $(*)$.
\underline{Then the pair}

\begin{equation*}
 \bigl(\underline R^{\mathrm{ab}},
 \underline F^{\mathrm{ab}}\subset
 \underline R^{\mathrm{ab}}\otimes_{\mathbb Z}
 \underline{\mathcal O}_S=\mathcal E^{\mathrm{ab}}\bigr)
\end{equation*}

\noindent
\underline{is isomorphic to the pair}
$(\underline R(B),\underline F(B))$
\underline{associated with an abeloid group over $S$}, as explained
after (13.4.4).  Thus

\begin{equation*}
 B\simeq\mathcal E^{\mathrm{ab}}/
 (F^{\mathrm{ab}}+R^{\mathrm{ab}}),
\end{equation*}

\noindent
where the symbols without underlining denote the corresponding analytic
groups over $S$.

\medskip
\noindent
\underline{Remark 13.4.18.} In Deligne's terminology [4], the conclusion
of the preceding proposition, after restriction to an open subset $U$
of $S$ on which $R$ is locally constant, says that $R$, endowed with the
two filtrations (13.4.17.1) on
$\mathcal E=\underline R\otimes_{\mathbb Z}\underline{\mathcal O}_S$,
is an ``analytic family of mixed Hodge structures over $U$''.  The
situation just described served as the model for Deligne's conjectures
on N\'eron theory in higher dimensions, stated in [9, 9.8--9.13].

There is, however, an important feature of the algebraic theory studied
in \S~7 which is missing from the description in 13.4.17.  Suppose that
$A$ is an abeloid group over an open subset $U$ complementary to a
nowhere-dense analytic subset $Z$ of $S$---say $Z=\{s\}$ when
$\dim S=1$---so that, by (13.1.5), $R|U$ is a local system over $U$.
Then the local system $R/R^t$ over $U$ ought to be \underline{constant};
more precisely, it ought to be dual to the local system ${R'}^f$, where
$R'$ is the analytic group, \`etale over $S$, which would correspond
(as $R$ corresponds to $A$) to a suitable analytic group $A'$ over $S$
such that $A'|U$ is the abeloid group dual to $A|U$.  In the notation of
\S~5, the algebraic analogue is the duality between
$T_\ell(A_K)/T_\ell(A_K)^t$ and $T_\ell(A'_K)^f$ in (5.5.9), which
follows from the duality theorem 5.2.

I do not know whether such a statement holds under only the hypotheses
of 13.4.17, or even whether the action of $\pi_1(U,t)$ on $R_t$ is
unipotent for $t\in U$.  Nor do the naive methods used in the present
paragraph seem capable of proving such a result, even under the
additional hypothesis that $A|U$ is endowed with a ``polarization'',
that is, with an invertible sheaf ample on every fibre.  This is,
however, something that the ``semi-simple group people'' apparently know
how to prove, at least when $S$ is nonsingular of dimension $1$; see the
report [9].  In the next subsection we shall show that everything works
under relative algebraicity hypotheses on $A$, and at the same time
relate the result to the construction of \S~7.
